\chapter{Inverse Preference Estimation Extension}

\section{Motivation}

Forward GOOP planning assumes that objective hierarchies are known. In realistic
traffic-like interaction, this is not true: other agents typically do not expose
their preference ordering. This chapter introduces the current estimation
extension built in \texttt{InverseGoopSolver}.

\section{Repository workflow}

The extension is organized into three executable programs:

\begin{description}
    \item[Program 1] centralized forward baseline,
    \item[Program 2] decentralized per-agent solving with prediction beliefs,
    \item[Program 3] ego preference-order identification by candidate testing.
\end{description}

\section{Program 3 estimation logic}

Program 3 currently estimates ordering by permutation search over a candidate
set for the ego agent:

\begin{enumerate}
    \item generate an oracle rollout with known ordering,
    \item test candidate orderings over an identification horizon,
    \item compute mismatch score between predicted and observed controls,
    \item select minimum-score ordering,
    \item run decentralized simulation with inferred order.
\end{enumerate}

Current score:
\begin{equation}
S(\pi) = \sum_{t=1}^{T_{\mathrm{id}}}
\|u^{\mathrm{pred}}_{\pi,t} - u^{\mathrm{obs}}_t\|_2^2.
\label{eq:score_inverse_chapter}
\end{equation}

\begin{figure}[tbp]
    \centering
    \fbox{\parbox{0.92\linewidth}{
    \textbf{Placeholder figure: Inverse estimation pipeline.}
    Show oracle rollout, candidate-order loop, score computation
    (Eq.~\ref{eq:score_inverse_chapter}), best-order selection, and
    decentralized rerun.
    }}
    \caption{Current inverse preference-order estimation workflow.}
    \label{fig:inverse_pipeline}
\end{figure}

\section{Relation to ongoing manuscript}

The ongoing RA-L manuscript direction uses this implementation as an executable
baseline to study preference-relation inference under partial observations. The
thesis keeps this work explicit as \emph{in-progress research}, not as finished
final results.

\section{Current limitations}

Main limitations at this stage:

\begin{itemize}
    \item factorial candidate growth with objective count,
    \item limited robustness analysis under noisy/partial observations,
    \item partial integration with embodied Duckietown experiments.
\end{itemize}

\section{Planned upgrades}

Planned improvements for final thesis stage include candidate pruning,
uncertainty-aware scoring, and broader benchmark coverage
\cite{peters:2023,li:2023,liu:2024}.
